% 原书分章 PDF 第 16—17 页：引理 23.1 的完整证明与本章末尾。
\begin{proof}[证明（引理 23.1）]
我们将证明 $\eta_n\to\eta_0$ 依概率成立，再由连续映射定理得到 $\tau_n$ 的一致性。由 23.3，存在一个包含 $\eta_0$ 的开区间 $\mathcal N$，使 $S_0$ 在其上连续。固定 $\eta\in\mathcal N$。由于 $\bar Q_{b,n}$ 以趋于 1 的概率属于一个 Glivenko--Cantelli 类，有
\begin{equation}
\begin{aligned}
 |S_n(\eta)-S_0(\eta)|
 &=|P_nI(\bar Q_{b,n}>\eta)-P_0I(\bar Q_{b,0}>\eta)|\\
 &\leq\left|P_0\bigl(I(\bar Q_{b,n}>\eta)-I(\bar Q_{b,0}>\eta)\bigr)\right|\\
 &\quad+|(P_n-P_0)I(\bar Q_{b,n}>\eta)|\\
 &\leq\underbrace{\left|P_0\bigl(I(\bar Q_{b,n}>\eta)-I(\bar Q_{b,0}>\eta)\bigr)\right|}_{\triangleq T_n(\eta)}
 +o_{P_0}(1),
\end{aligned}
\tag{23.12}
\end{equation}
\Needspace{4\baselineskip}
其中，对任意分布 $P$ 与函数 $f$，采用记号 $Pf=\E_P[f(O)]$。令
$Z_n(\eta)(w)\triangleq\bigl(I(\bar Q_{b,n}(w)>\eta)-I(\bar Q_{b,0}(w)>\eta)\bigr)^2$。
对所有 $q>0$，有下列各式：
\begin{align}
 T_n(\eta)&\leq P_0Z_n(\eta)\nonumber\\
 &=P_0Z_n(\eta)I(|\bar Q_{b,0}-\eta|>q)
 +P_0Z_n(\eta)I(|\bar Q_{b,0}-\eta|\leq q)\nonumber\\
 &=P_0Z_n(\eta)I(|\bar Q_{b,0}-\eta|>q)
 +P_0Z_n(\eta)I(0<|\bar Q_{b,0}-\eta|\leq q)\tag{23.13}\\
 &\leq P_0Z_n(\eta)I(|\bar Q_{b,n}-\bar Q_{b,0}|>q)
 +P_0Z_n(\eta)I(0<|\bar Q_{b,0}-\eta|\leq q)\tag{23.14}\\
 &\leq P_0(|\bar Q_{b,n}-\bar Q_{b,0}|>q)
 +P_0(0<|\bar Q_{b,0}-\eta|\leq q)\nonumber\\
 &\leq\frac{P_0|\bar Q_{b,n}-\bar Q_{b,0}|}{q}
 +P_0(0<|\bar Q_{b,0}-\eta|\leq q).\nonumber
\end{align}
上式中，式 (23.13) 成立是因为 23.3 蕴含 $P_0(\bar Q_{b,0}=\eta)=0$；式 (23.14) 成立是因为 $Z_n(\eta)=1$ 蕴含
$|\bar Q_{b,n}-\bar Q_{b,0}|\geq|\bar Q_{b,0}-\eta|$；最后一个不等式由 Markov 不等式成立。引理假定 $\E_0|\bar Q_{b,n}-\bar Q_{b,0}|=o_{P_0}(1)$，因此可以选择一个序列 $q_n\downarrow0$，使得
\[
 T_n(\eta)\leq P_0(0<|\bar Q_{b,0}-\eta|\leq q_n)+o_{P_0}(1).
\]
为说明右侧第一项为 $o(1)$，注意，$P_0(\bar Q_{b,0}=\eta)=0$ 与 $S_0$ 在 $\mathcal N$ 上的连续性共同给出：对于足够大的 $n$，
\[
 P_0(0<|\bar Q_{b,0}-\eta|\leq q_n)=S_0(-q_n+\eta)-S_0(q_n+\eta).
\]
右侧为 $o(1)$，因此 $T_n(\eta)=o_{P_0}(1)$。代入式 (23.12)，可知 $S_n(\eta)\to S_0(\eta)$ 依概率成立。回顾 $\eta\in\mathcal N$ 是任意的。

固定 $\gamma>0$。当 $\gamma$ 足够小时，$\eta_0-\gamma$ 和 $\eta_0+\gamma$ 均包含于 $\mathcal N$。于是，
$S_n(\eta_0-\gamma)\to S_0(\eta_0-\gamma)$ 和 $S_n(\eta_0+\gamma)\to S_0(\eta_0+\gamma)$ 依概率成立。
此外，由 $\eta_0$ 的定义，$S_0(\eta_0-\gamma)>\kappa$；由条件 23.3，$S_0(\eta_0+\gamma)<\kappa$。
因此，以趋于 1 的概率，有 $S_n(\eta_0-\gamma)>\kappa$ 且 $S_n(\eta_0+\gamma)<\kappa$。
但是，$|\eta_n-\eta_0|>\gamma$ 蕴含 $S_n(\eta_0-\gamma)\leq\kappa$ 或 $S_n(\eta_0+\gamma)>\kappa$，所以以趋于 1 的概率有 $|\eta_n-\eta_0|\leq\gamma$。
于是 $\eta_n\to\eta_0$ 依概率成立；由连续映射定理，$\tau_n\to\tau_0$。
\end{proof}
